% ══════════════════════════════════════════════════════════════════════════════
% Supplementary Information
% ══════════════════════════════════════════════════════════════════════════════

% Reset counters for supplementary numbering
\setcounter{section}{0}
\setcounter{figure}{0}
\setcounter{table}{0}
\setcounter{equation}{0}
\renewcommand{\thesection}{S\arabic{section}}
\renewcommand{\thefigure}{S\arabic{figure}}
\renewcommand{\thetable}{S\arabic{table}}
\renewcommand{\theequation}{S\arabic{equation}}

\section*{Supplementary Information}

% ══════════════════════════════════════════════════════════════════════════════
\subsection*{Supplementary Tables}

% ── Table S1: Beta sweep ─────────────────────────────────────────────────
\begin{table}[ht]
\centering
\caption{\textbf{Inverse temperature sensitivity analysis.} SA generation
quality as a function of $\beta / \beta^*$, where $\beta^* \approx 2.94$ was
identified via entropy inflection in the 18-dimensional PCA space ($K{=}23$
stored patterns). Operating at $\beta = \beta^*$ (bold row) balances generation
novelty against fidelity. Med MRE = median relative error of per-feature means;
Novelty = $1 - \max_k \cos(\hat{\xi}, \hat{m}_k)$.}
\label{stab:beta-sweep}
\begin{tabular}{rrcccc}
\toprule
$\beta/\beta^*$ & $\beta$ & Noise scale & PC1 std ratio & Novelty & Med MRE (\%) \\
\midrule
0.1 & 0.29 & 0.261 & 0.538 & 0.542 & 0.7 \\
0.3 & 0.88 & 0.151 & 0.550 & 0.536 & 0.7 \\
0.5 & 1.47 & 0.117 & 0.562 & 0.528 & 0.7 \\
\textbf{1.0} & \textbf{2.94} & \textbf{0.082} & \textbf{0.589} & \textbf{0.502} & \textbf{0.6} \\
1.5 & 4.41 & 0.067 & 0.613 & 0.474 & 0.7 \\
2.0 & 5.88 & 0.058 & 0.638 & 0.446 & 0.7 \\
3.0 & 8.82 & 0.048 & 0.689 & 0.403 & 0.8 \\
\bottomrule
\end{tabular}
\end{table}

% ── Table S2: Multiplicity weighting parameters ─────────────────────────
\begin{table}[ht]
\centering
\caption{\textbf{Multiplicity weighting parameters for conditional generation.}
The multiplicity ratio $\rho$ was computed to achieve a target effective fraction
$f_{\text{target}} = 0.80$ of softmax attention on the designated patterns.
$K_{\text{eff}}$ is the participation ratio quantifying the effective number of
patterns contributing to generation.}
\label{stab:multiplicity}
\begin{tabular}{lccccc}
\toprule
Condition & $K_{\text{sub}}$ & $K_{\text{bg}}$ & $\rho$ & $f_{\text{eff}}$ & $K_{\text{eff}}$ \\
\midrule
Uncomplicated  & 18 & 5  & 1.11  & 0.80 & 23.0 \\
PCOS           & 3  & 20 & 26.67 & 0.80 & 4.6 \\
Developed PE   & 5  & 18 & 14.40 & 0.80 & 7.7 \\
\bottomrule
\end{tabular}
\end{table}

% ── Table S3: SA Hyperparameters ─────────────────────────────────────────
\begin{table}[ht]
\centering
\caption{\textbf{SA generation hyperparameters.} The same values were used for
unconditioned and conditioned generation except where noted.}
\label{stab:hyperparams}
\begin{tabular}{lll}
\toprule
Parameter & Value & Description \\
\midrule
$\alpha$ & 0.01 & Langevin step size (ULA valid without MH correction) \\
$T$ & 2,000 & Langevin iterations per sample \\
$\beta^*$ & 2.94 & Inverse temperature (entropy inflection) \\
PCA threshold & 95\% & Variance retained \\
$d_{\text{PCA}}$ & 18 & PCA dimensionality \\
$N$ & 100 & Synthetic patients per cohort \\
Random seed & 42 & For reproducibility \\
$f_{\text{target}}$ & 0.80 & Target attention fraction (conditioned only) \\
\bottomrule
\end{tabular}
\end{table}

% ── Table S4: ODE Calibration Details ────────────────────────────────────
\begin{table}[ht]
\centering
\caption{\textbf{BZ2012 ODE model calibration.} Five rate constants were
calibrated on Visit~1 real patients ($n{=}23$) using bounded Nelder--Mead
optimization in log-scale-factor space. The cost function was the mean
normalized relative squared error across five TGA features under TF-only
conditions. Bounds: $\pm\log(100)$. Convergence: $f_{\text{tol}} = 10^{-6}$,
500 iterations, 11 random restarts.}
\label{stab:ode-calibration}
\begin{tabular}{llrrl}
\toprule
Rate constant & Role & Default & Calibrated & Scale \\
\midrule
prothrombinase $k_{\text{cat}}$ & Thrombin burst & 63.5 s$^{-1}$ & 21.94 s$^{-1}$ & 0.35$\times$ \\
intrinsic Xase $k_{\text{cat}}$ & Amplification Xa & 8.2 s$^{-1}$ & 0.169 s$^{-1}$ & 0.021$\times$ \\
extrinsic Xase $k_{\text{cat}}$ & Initiation Xa & 6.0 s$^{-1}$ & 93.2 s$^{-1}$ & 15.5$\times$ \\
PC activation $k_{\text{cat}}$ & Protein C activation & 0.41 s$^{-1}$ & 0.890 s$^{-1}$ & 2.17$\times$ \\
mIIa conversion $k$ & mIIa$\to$IIa & $2.3{\times}10^8$ & $1.15{\times}10^9$ & 5.01$\times$ \\
\bottomrule
\multicolumn{5}{l}{\small Units for mIIa conversion: M$^{-1}$s$^{-1}$} \\
\end{tabular}
\end{table}

% ── Table S5: Full MRE summary ───────────────────────────────────────────
\begin{table}[ht]
\centering
\caption{\textbf{Summary of mean relative errors across all 72 features.}
Per-visit distribution of MREs between real and SA-generated synthetic patient
means. The full feature-by-feature table (72 features $\times$ 3 visits = 216
entries) is provided in the code repository as
\texttt{paper\_summary\_statistics.csv}.}
\label{stab:mre-summary}
\begin{tabular}{lccccc}
\toprule
Visit & Median MRE & Mean MRE & 25th pctl & 75th pctl & Max MRE \\
\midrule
V1 (baseline) & 0.019 & 0.027 & 0.010 & 0.033 & 0.157 \\
V2 (1st tri)  & 0.019 & 0.027 & 0.009 & 0.034 & 0.124 \\
V3 (3rd tri)  & 0.016 & 0.025 & 0.006 & 0.040 & 0.110 \\
\bottomrule
\end{tabular}
\end{table}

% ── Table S6: CTGAN/TVAE comparison ──────────────────────────────────────
\begin{table}[ht]
\centering
\caption{\textbf{Comparison with deep generative baselines.} CTGAN and TVAE
were trained on the same 90 per-visit records (23 patients $\times$ 3 visits,
72 features) at three epoch counts. SA and MVN results from the main text are
shown for reference. CTGAN fails at this sample size (median MRE
$\approx$~19\%), consistent with GAN mode collapse on small datasets. TVAE
achieves comparable marginal fidelity at high epoch counts but cannot capture
cross-visit covariance or perform conditional generation.}
\label{stab:ctgan}
\begin{tabular}{lcccc}
\toprule
Method & Epochs & Med MRE & Med KS & Corr MAE \\
\midrule
SA              & --   & 0.019 & 0.169 & 0.200 \\
MVN             & --   & 0.024 & 0.161 & 0.090 \\
\midrule
CTGAN           & 300   & 0.189 & 0.364 & 0.261 \\
CTGAN           & 1,000 & 0.195 & 0.349 & 0.265 \\
CTGAN           & 3,000 & 0.187 & 0.341 & 0.180 \\
\midrule
TVAE            & 300   & 0.037 & 0.174 & 0.152 \\
TVAE            & 1,000 & 0.026 & 0.248 & 0.082 \\
TVAE            & 3,000 & 0.018 & 0.224 & 0.092 \\
\bottomrule
\end{tabular}
\end{table}

% ── Table S7: PCA sensitivity ────────────────────────────────────────────
\begin{table}[ht]
\centering
\caption{\textbf{PCA variance threshold sensitivity analysis.} SA generation
quality is stable across thresholds from 85\% to 99\%: median MRE remains
between 1.2\% and 1.8\%. The 95\% threshold used in the main text (bold) is
not a critical design choice.}
\label{stab:pca-sensitivity}
\begin{tabular}{rrccccc}
\toprule
Threshold & $d_{\text{PCA}}$ & $K/d$ & $\beta^*$ & Novelty & Med MRE & Corr MAE \\
\midrule
85\% & 13 & 1.77 & 2.94 & 0.413 & 0.012 & 0.116 \\
90\% & 15 & 1.53 & 2.94 & 0.462 & 0.018 & 0.131 \\
\textbf{95\%} & \textbf{18} & \textbf{1.28} & \textbf{2.94} & \textbf{0.516} & \textbf{0.012} & \textbf{0.148} \\
98\% & 20 & 1.15 & 2.94 & 0.538 & 0.013 & 0.127 \\
99\% & 21 & 1.10 & 2.94 & 0.540 & 0.013 & 0.133 \\
\bottomrule
\end{tabular}
\end{table}

\FloatBarrier

% ══════════════════════════════════════════════════════════════════════════════
\subsection*{Supplementary Figures}

% ── Fig S1: Cross-visit corr amplified ───────────────────────────────────
\begin{figure}[ht]
\centering
\includegraphics[width=\textwidth]{validate_cross_visit_corr_supp.pdf}
\caption{\textbf{Cross-visit correlation structure (amplified residual scale).}
Same layout as the main-text correlation figure, but with the bottom-row
residual panels plotted on a $\pm 0.5$ color scale (vs.\ $\pm 1.0$ in the main
text) to reveal fine structure. Frobenius norms of each residual matrix are
shown in the panel titles.}
\label{sfig:cross-visit-corr-supp}
\end{figure}

% ── Fig S2: Eigenvalue spectrum ──────────────────────────────────────────
\begin{figure}[ht]
\centering
\includegraphics[width=0.75\textwidth]{validate_eigenvalue_spectrum.pdf}
\caption{\textbf{Covariance eigenvalue spectrum.} Log-scale eigenvalues of the
$216 \times 216$ sample covariance for Real ($K{=}23$, rank~22), SA ($N{=}100$,
rank~18), and MVN ($N{=}100$) populations. Dashed vertical lines mark the SA
PCA truncation point (component~18) and the data rank limit (component~22).
MVN's Ledoit--Wolf regularization inflates eigenvalues beyond component~22
(green shaded region), introducing spurious variance in 194 dimensions where the
data contain no signal.}
\label{sfig:eigenvalue-spectrum}
\end{figure}

% ── Fig S3: Conditioned generation PCA (pooled) ─────────────────────────
\begin{figure}[ht]
\centering
\includegraphics[width=0.65\textwidth]{validate_conditioned_pca.pdf}
\caption{\textbf{Conditioned generation: PCA projection (pooled).} Large
markers: real patients (blue = Uncomplicated, $n{=}18$; orange = PCOS, $n{=}3$;
red = Developed~PE, $n{=}5$). Small markers: synthetic patients ($N{=}100$ per
condition). With only 3 PCOS and 5 Developed~PE patients, the subgroups do not form
clearly separable clusters in the first two principal components, but the
synthetic cohorts concentrate around their respective real counterparts.}
\label{sfig:conditioned-pca}
\end{figure}

% ── Fig S4: Conditioned generation by visit ──────────────────────────────
\begin{figure}[ht]
\centering
\includegraphics[width=\textwidth]{validate_conditioned_pca_by_visit.pdf}
\caption{\textbf{Conditioned generation: PCA projection by visit.} Same as
Supplementary Fig.~\ref{sfig:conditioned-pca} but separated by visit (V1, V2,
V3). Condition-specific clustering is maintained within each individual visit.}
\label{sfig:conditioned-by-visit}
\end{figure}

% ── Fig S5: Mechanistic TF+TM (combined scatter + ratios) ───────────────
\begin{figure}[ht]
\centering
\includegraphics[width=\textwidth]{validate_mechanistic_combined_tf_tm.pdf}
\caption{\textbf{Mechanistic validation (TF+TM condition).} Same layout as
main-text Fig.~6 but under TF+TM conditions. The model systematically
underpredicts peak and maximum rate ($\approx 0.5\times$ measured) for both
real and synthetic patients, reflecting overcorrection of the protein~C
anticoagulant pathway. Cloud overlap remains high at 88--92\%, confirming that
the model processes real and synthetic patients identically even under imperfect
calibration conditions.}
\label{sfig:mechanistic-tm-combined}
\end{figure}

% ── Fig S6: Cross-visit generalization (TF-only + TF+TM) ────────────────
\begin{figure}[ht]
\centering
\includegraphics[width=\textwidth]{validate_mechanistic_generalization_tf_only.pdf}\\[6pt]
\includegraphics[width=\textwidth]{validate_mechanistic_generalization_tf_tm.pdf}
\caption{\textbf{Cross-visit calibration generalization.} Boxplots of the
ODE-predicted/dataset ratio by visit for real patients. The model was calibrated
on Visit~1 only. Top: TF-only conditions, where ratios near 1.0 at Visits~2
and~3 indicate that the calibration generalizes across pregnancy timepoints.
Bottom: TF+TM conditions, where the calibration does not generalize as well,
particularly for peak and maximum rate.}
\label{sfig:generalization}
\end{figure}

% ── Fig S7: Rank correlations (TF-only + TF+TM) ─────────────────────────
\begin{figure}[ht]
\centering
\includegraphics[width=0.48\textwidth]{validate_mechanistic_rankcorr_tf_only.pdf}
\hfill
\includegraphics[width=0.48\textwidth]{validate_mechanistic_rankcorr_tf_tm.pdf}
\caption{\textbf{Spearman rank correlations between ODE-predicted and dataset
TGA features.} Left: TF-only. Right: TF+TM. ETP is the best-predicted feature
($\rho \approx 0.6$--$0.8$ for real patients under TF-only). Weak rank
correlations for other features are expected given the 5-parameter
population-level calibration.}
\label{sfig:rankcorr}
\end{figure}
